\documentclass[11pt]{article}
\usepackage{mathblatt}
\begin{document}
% Kopfzeile Seite 1 ohne Kennung; die Rückseite (Lösungen, für den Lehrer) bekommt sie nach \begleitteil
\blattkopf{Mathematik}{Basisaufgaben}
\einheitenkopf[][Zettel 16]{Basisaufgaben · Zettel 16}
\zweigzeile{je Aufgabe 1 Punkt · Lösungen auf der Rückseite}
\par\vspace{2pt}\noindent Name \underline{\hspace{7cm}}\hfill Datum \underline{\hspace{3cm}}\hfill Punkte \underline{\hspace{1.2cm}} / 17\par\vspace{6pt}
% \small: volle Seite, kurze Aufgaben zweispaltig (v0.6)
\small

\noindent\begin{minipage}[t]{0.485\linewidth}\raggedright
% 1: Bruchteil einer Größe berechnen – brueche-dezimalzahlen-basis-k2-v6 (Original 2018-OS-B1a, ausgedacht: keine Marke)
\begin{aufgabe}{Wie viel Meter sind $\frac{3}{4}$ von $1{,}6$ km? \leerfeld[m]}
\end{aufgabe}
\end{minipage}\hfill\begin{minipage}[t]{0.485\linewidth}\raggedright
% 2: Zahl zu Bedingung angeben – rationale-zahlen-basis-k3-v4 (Original 2015-OS-B1b, ausgedacht: keine Marke)
\begin{aufgabe}{Gib eine Zahl an, die zwischen $\frac{1}{4}$ und $\frac{1}{2}$ liegt. \leerfeld}
\end{aufgabe}
\end{minipage}\par

\noindent\begin{minipage}[t]{0.485\linewidth}\raggedright
% 3: Exponent einer Potenz bestimmen – potenzen-wurzeln-basis-k1-v6 (Original 2024-OS-B1g, ausgedacht: keine Marke)
\begin{aufgabe}{$3^x = 27$ – welche Zahl ist $x$? x = \leerfeld}
\end{aufgabe}
\end{minipage}\hfill\begin{minipage}[t]{0.485\linewidth}\raggedright
% 4: Zehnerpotenzschreibweise umwandeln – potenzen-wurzeln-basis-k3-v8 (Original 2025-OS-B1d, ausgedacht: keine Marke)
\begin{aufgabe}{$0{,}00036 = 3{,}6 \cdot 10^{\square}$ – welche Hochzahl? \leerfeld}
\end{aufgabe}
\end{minipage}\par

\noindent\begin{minipage}[t]{0.485\linewidth}\raggedright
% 5: Wert nach prozentualer Erhöhung berechnen – prozentrechnung-basis-k5-v2 (Original 2024-OS-B1e, ausgedacht: keine Marke)
\begin{aufgabe}{Ein Buch kostet $8\,€$ und wird um $50\,\%$ teurer. Wie hoch ist der neue Preis? \leerfeld[€]}
\end{aufgabe}
\end{minipage}\hfill\begin{minipage}[t]{0.485\linewidth}\raggedright
% 6: Zeiteinheiten umrechnen – einheiten-basis-k5-v3 (Original 2025-OS-B1b, ausgedacht: keine Marke)
\begin{aufgabe}{Eine Wanderung dauert $4{,}25$ h. Gib die Dauer in Minuten an. \leerfeld[min]}
\end{aufgabe}
\end{minipage}\par

\noindent\begin{minipage}[t]{0.485\linewidth}\raggedright
% 7: Lineare Gleichung lösen – lineare-gleichungen-basis-k1-v10 (Original 2024-OS-B1d, ausgedacht: keine Marke)
\begin{aufgabe}{Löse die Gleichung $4 \cdot (x - 1) + 2 = 10$. x = \leerfeld}
\end{aufgabe}
\end{minipage}\hfill\begin{minipage}[t]{0.485\linewidth}\raggedright
% 8: Arithmetisches Mittel berechnen – daten-basis-k1-v4 (Original 2021-OS-B1f, ausgedacht: keine Marke)
\begin{aufgabe}{Ein Heft kostet in vier Läden $2{,}40$; $3{,}10$; $2{,}80$; $3{,}30$ €. Berechne das arithmetische Mittel (Preis). \leerfeld[€]}
\end{aufgabe}
\end{minipage}\par

% 9: Wurzel eines Quadrats berechnen – potenzen-wurzeln-basis-k2-v2 (Original 2015-OS-B1j, ausgedacht: keine Marke)
\begin{aufgabe}{Welche Aussage ist wahr? Kreuze an.\\ \kreuz{$\sqrt{(-9)^2} = -9$}\\ \kreuz{$\sqrt{(-9)^2} = 9$}\\ \kreuz{$\sqrt{(-9)^2}$ ist nicht definiert}}
\end{aufgabe}

% 10: Antiproportionale Zuordnung Dreisatz – zuordnungen-basis-k1-v4 (Original 2014-GYM-B1c, ausgedacht: keine Marke)
\begin{aufgabe}{Für einen Ausflug braucht man 3 Busse mit je 40 Plätzen. Wie viele Plätze müsste jeder Bus haben, wenn nur 2 Busse fahren? \leerfeld Plätze}
\end{aufgabe}

% 11: Term zu Sachtext angeben – terme-basis-k2-v9 (Original 2023-OS-B1h, ausgedacht: keine Marke)
\begin{aufgabe}{Die Differenz aus dem Siebenfachen einer Zahl $y$ und 2 wird vervierfacht. Welcher Term passt? \\ \kreuz{$(7y - 2) : 4$} \kreuz{$7y - 2 \cdot 4$} \kreuz{$4 \cdot (7y - 2)$} \kreuz{$7y \cdot 4 - 2$}}
\end{aufgabe}

% 12: Dreiecksungleichung anwenden – winkel-dreiecke-basis-k1-v4 (Original 2020-OS-B1i, ausgedacht: keine Marke)
\begin{aufgabe}{Aus welchen drei Stäben kann man ein Dreieck legen? \\ \kreuz{$5$ cm, $5$ cm, $10$ cm} \kreuz{$2$ cm, $7$ cm, $4$ cm} \kreuz{$6$ cm, $3$ cm, $8$ cm}}
\end{aufgabe}

% 13: Eigenschaft einer Figur zuordnen – winkel-dreiecke-basis-k5-v4 (Original 2026-FOR-B1c, ausgedacht: keine Marke)
\begin{aufgabe}{Kreuze die richtige Ergänzung an: „In jedem Drachenviereck …“\\ \kreuz{… stehen die Diagonalen senkrecht aufeinander.}\\ \kreuz{… sind gegenüberliegende Seiten parallel.}\\ \kreuz{… sind alle Seiten gleich lang.}}
\end{aufgabe}

% 14: Ergebnismenge aufzählen – wahrscheinlichkeit-basis-k1-v2 (Original 2017-OS-B1f, ausgedacht: keine Marke)
\begin{aufgabe}{Ein Glücksrad mit den Farben Rot, Gelb und Blau wird zweimal gedreht. Wie viele verschiedene Ergebnisse gibt es? \leerfeld Ergebnisse}
\end{aufgabe}

% 15: Zufallsgerät zu Wahrscheinlichkeit entwerfen – wahrscheinlichkeit-basis-k2-v3 (Original 2019-OS-B1g, ausgedacht: keine Marke)
\begin{aufgabe}{Ein Stapel hat $12$ Karten. Die Wahrscheinlichkeit, einen Joker zu ziehen, soll $\frac{1}{4}$ sein. Wie viele Joker braucht man? \leerfeld Joker}
\end{aufgabe}

% 16: Symmetrieachsen bestimmen – symmetrie-abbildungen-basis-k3-v3 (Original 2022-OS-B1i, ausgedacht: keine Marke)
\begin{aufgabe}{\begin{minipage}[t]{0.6\linewidth}Ein Drachen hat die Form eines Drachenvierecks (keine Raute). Wie viele Symmetrieachsen hat die Figur? Kreuze an. \\ \kreuz{0} \kreuz{1} \kreuz{2} \kreuz{3} \kreuz{4} \kreuz{5}\end{minipage}\hfill\begin{minipage}[t]{0.37\linewidth}\vspace{0pt}
\drachen{2.4}{1.8}{0.3}
\end{minipage}}
\end{aufgabe}

% 17: Kreissektor Anteil berechnen – kreis-basis-k1-v2 (Original 2025-OS-B1e, ausgedacht: keine Marke)
\begin{aufgabe}{\begin{minipage}[t]{0.6\linewidth}Auf dem grauen Teil eines runden Beets wachsen Rosen. Der Mittelpunktswinkel des grauen Teils beträgt $72^\circ$. Wie viel Prozent der Kreisfläche sind grau? \leerfeld[\%]\end{minipage}\hfill\begin{minipage}[t]{0.37\linewidth}\vspace{0pt}
\kreissektor[1.2]{72}{$72^\circ$}
\end{minipage}}
\end{aufgabe}

\begleitteil
\blattkopf{Basisaufgaben · Zettel 16}{BAS-Z16 · Lösungen}
\erg{1}{$1{,}6$ km $= 1\,600$ m, $1\,600 : 4 \cdot 3 = 1\,200$ m \hfill {\footnotesize\textit{Bruchteil einer Größe berechnen · 2018}}}
\erg{2}{z. B. $0{,}3$; richtig ist jede Zahl zwischen $0{,}25$ und $0{,}5$ \hfill {\footnotesize\textit{Zahl zu Bedingung angeben · 2015}}}
\erg{3}{$3, 9, 27$: drei Faktoren, also $x = 3$ \hfill {\footnotesize\textit{Exponent einer Potenz bestimmen · 2024}}}
\erg{4}{$-4$ – das Komma wandert vier Stellen nach rechts \hfill {\footnotesize\textit{Zehnerpotenzschreibweise umwandeln · 2025}}}
\erg{5}{$8 \cdot 1{,}5 = 12\,€$ (nicht nur $4\,€$) \hfill {\footnotesize\textit{Wert nach prozentualer Erhöhung berechnen · 2024}}}
\erg{6}{$4{,}25 \cdot 60 = 255$ min \hfill {\footnotesize\textit{Zeiteinheiten umrechnen · 2025}}}
\erg{7}{$4 \cdot (x - 1) = 8$, $x - 1 = 2$, also $x = 3$ \hfill {\footnotesize\textit{Lineare Gleichung lösen · 2024}}}
\erg{8}{$11{,}60 : 4 = 2{,}90$ € \hfill {\footnotesize\textit{Arithmetisches Mittel berechnen · 2021}}}
\erg{9}{$\sqrt{(-9)^2} = 9$, denn erst das Quadrat: $(-9)^2 = 81$, und die Wurzel ist nie negativ \hfill {\footnotesize\textit{Wurzel eines Quadrats berechnen · 2015}}}
\erg{10}{$3 \cdot 40 = 120$ Plätze; $120 : 2 = 60$ Plätze \hfill {\footnotesize\textit{Antiproportionale Zuordnung Dreisatz · 2014}}}
\erg{11}{$4 \cdot (7y - 2)$ \hfill {\footnotesize\textit{Term zu Sachtext angeben · 2023}}}
\erg{12}{$6$ cm, $3$ cm, $8$ cm \hfill {\footnotesize\textit{Dreiecksungleichung anwenden · 2020}}}
\erg{13}{… stehen die Diagonalen senkrecht aufeinander. \hfill {\footnotesize\textit{Eigenschaft einer Figur zuordnen · 2026}}}
\erg{14}{$3 \cdot 3 = 9$ (RR, RG, RB, GR, GG, GB, BR, BG, BB) \hfill {\footnotesize\textit{Ergebnismenge aufzählen · 2017}}}
\erg{15}{$12 : 4 = 3$ \hfill {\footnotesize\textit{Zufallsgerät zu Wahrscheinlichkeit entwerfen · 2019}}}
\erg{16}{1 \hfill {\footnotesize\textit{Symmetrieachsen bestimmen · 2022}}}
\erg{17}{$72 : 360 = 0{,}2 = 20\,\%$ \hfill {\footnotesize\textit{Kreissektor Anteil berechnen · 2025}}}
\end{document}
